\section*{Chapitre \ref{ch:b2:huckel} --- Théorie de Hückel et systèmes conjugués}

\begin{solution}{exo:b2:huckel:1}
Niveaux de l'allyle ($n = 3$)~: $\alpha + 1.414\beta$, $\alpha$, $\alpha - 1.414\beta$. Cation (2 électrons)~:
$2\alpha + 2.828\beta$~; radical (3)~: $3\alpha + 2.828\beta$~; anion (4)~: $4\alpha + 2.828\beta$. Les électrons
supplémentaires vont dans le niveau non liant et ne modifient que la partie en $\alpha$.
\end{solution}

\begin{solution}{exo:b2:huckel:2}
$2(\alpha + 1.618\beta) + 2(\alpha + 0.618\beta) = 4\alpha + 4.472\beta$.
\end{solution}

\begin{solution}{exo:b2:huckel:3}
Benzène (6), anion cyclopentadiényle (6), tropylium (6)~: \omterm{def:b2:huckel:aromatic}{aromatiques}. Cation
cyclopentadiényle (4), cyclobutadiène (4), cyclooctatétraène plan (8)~: \omterm{def:b2:huckel:aromatic}{antiaromatiques} s'ils sont plans
(le cyclooctatétraène réel a la forme d'une baignoire et n'est pas \omterm{def:b2:huckel:aromatic}{aromatique}).
\end{solution}

\begin{solution}{exo:b2:huckel:4}
Un octogone dans le cercle~: niveaux $\alpha + 2\beta$, $\alpha + 1.414\beta$ (deux fois), $\alpha$ (deux fois),
$\alpha - 1.414\beta$ (deux fois), $\alpha - 2\beta$. Huit électrons~: deux dans le plus bas, quatre dans la
première paire, un dans chaque orbitale de la paire non liante.
\end{solution}

\begin{solution}{exo:b2:huckel:5}
$4.472\beta - 4\beta = 0.472\beta$, soit $0.472 \times 75 = \qty{35}{kJ/mol}$, contre environ
\qty{10}{kJ/mol} d'après les données d'hydrogénation~: le modèle de Hückel, avec $\beta$ ajusté sur le benzène,
surestime la conjugaison des chaînes ouvertes.
\end{solution}

\begin{solution}{exo:b2:huckel:6}
$p_{12} = p_{34} = 0.894$, $p_{23} = 0.447$~: les liaisons terminales C1--C2 et C3--C4 sont les
plus courtes, la liaison centrale plus longue, mais plus courte qu'une liaison simple.
\end{solution}

\begin{solution}{exo:b2:huckel:7}
Anion, six électrons~: $2 \times 2 + 4 \times 0.618$, $E_\pi = 6\alpha + 6.472\beta$, une couche complète~:
\omterm{def:b2:huckel:aromatic}{aromatique}. Cation, quatre~: $2 \times 2 + 2 \times 0.618$, $4\alpha + 5.236\beta$, avec un électron
dans chaque orbitale de la paire dégénérée~: \omterm{def:b2:huckel:aromatic}{antiaromatique}, très instable.
\end{solution}

\begin{solution}{exo:b2:huckel:8}
Cycle à sept atomes~: $\alpha + 2\beta$, $\alpha + 1.247\beta$ (deux fois), puis des niveaux antiliants. Six
électrons, comme dans le cation \ce{C7H7+}, remplissent exactement les niveaux liants~: un cation
\omterm{def:b2:huckel:aromatic}{aromatique}. Le composé s'ionise facilement en tropylium et bromure.
\end{solution}

\begin{solution}{exo:b2:huckel:9}
$1.802 + 1.247 + 0.445 - 0.445 - 1.247 - 1.802 = 0$~: la somme des énergies vaut $6\alpha$.
\end{solution}

\begin{solution}{exo:b2:huckel:10}
$n = 6$~: $m_k = 2\cos(k\pi/7) = 1.802$, $1.247$, $0.445$, $-0.445$, $-1.247$, $-1.802$. Six
électrons~: $E_\pi = 6\alpha + 2(1.802 + 1.247 + 0.445)\beta = 6\alpha + 6.988\beta$~;
délocalisation $0.988\beta$.
\end{solution}

\begin{solution}{exo:b2:huckel:11}
Dans le carré, les deux électrons de plus haute énergie se placent chacun dans une orbitale d'une
paire non liante dégénérée. Allonger deux liaisons opposées lève la dégénérescence~:
une orbitale descend, l'autre monte, et les deux électrons s'apparient dans la plus basse, ce qui
abaisse l'énergie. La molécule se stabilise en rectangle avec deux doubles liaisons
localisées~; elle reste \omterm{def:b2:huckel:aromatic}{antiaromatique} et extrêmement réactive.
\end{solution}

\begin{solution}{exo:b2:huckel:12}
Matrice (en unités de $\beta$, avec $\alpha$ pour origine) $\begin{pmatrix}0 & 1\\ 1 & 1\end{pmatrix}$~:
$m = 1.618$ et $-0.618$, $E = \alpha + 1.618\beta$ (liante) et $\alpha - 0.618\beta$. Pour
l'\omterm{def:b2:diatomic-mos:bonding}{orbitale liante}, $c_X = 1.618c_C$, donc $c_C^2 = 0.276$, $c_X^2 = 0.724$~; avec deux électrons,
$q_C = 0.55$, $q_X = 1.45$~: la densité $\pi$ est attirée vers X.
\end{solution}

\begin{solution}{pb:b2:huckel:1}
\textbf{1.} $-123.1 - (-4.3) = \qty{-118.8}{kJ/mol}$.
\textbf{2.} $-123.1 - 104.6 = \qty{-227.7}{kJ/mol}$.
\textbf{3.} $-123.1 - 82.9 = \qty{-206.0}{kJ/mol}$.
\textbf{4.} $3 \times (-118.8) = \qty{-356.4}{kJ/mol}$.
\textbf{5.} $356.4 - 206.0 = \qty{150}{kJ/mol}$.
\textbf{6.} Le diène libère $2 \times 118.8 - 227.7 = \qty{10}{kJ/mol}$ de moins que deux
doubles liaisons isolées~: la conjugaison d'une chaîne ouverte apporte un faible gain, le cycle
fermé du benzène quinze fois plus.
\textbf{7.} Une matrice $6 \times 6$ portant 1 entre chaque atome et ses deux voisins (1--2, 2--3,
\dots, 6--1) et 0 ailleurs.
\textbf{8.} $m = 2\cos(2\pi k/6)$~: $2$, $1$, $1$, $-1$, $-1$, $-2$~: $E = \alpha + 2\beta$, $\alpha + \beta$ (deux fois),
$\alpha - \beta$ (deux fois), $\alpha - 2\beta$.
\textbf{9.} Deux électrons dans $\alpha + 2\beta$, quatre dans la paire $\alpha + \beta$.
\textbf{10.} $E_\pi = 6\alpha + 8\beta$.
\textbf{11.} $3(2\alpha + 2\beta) = 6\alpha + 6\beta$.
\textbf{12.} $2\beta$, soit $2|\beta|$ de stabilisation.
\textbf{13.} $2 + 1 + 1 - 1 - 1 - 2 = 0$~: les six niveaux ont pour somme $6\alpha$.
\textbf{14.} $2|\beta| = \qty{150}{kJ/mol}$~: $|\beta| = \qty{75}{kJ/mol}$.
\textbf{15.} $0.472 \times 75 = \qty{35}{kJ/mol}$, contre \qty{10}{kJ/mol}~: le modèle ajusté
sur le benzène surestime la chaîne ouverte.
\textbf{16.} On peut prendre pour orbitales occupées $\psi_1 = (1, 1, 1, 1, 1, 1)/\sqrt6$,
$\psi_2 = (2, 1, -1, -2, -1, 1)/\sqrt{12}$ et $\psi_3 = (0, 1, 1, 0, -1, -1)/2$. Pour la liaison
1--2~: $p_{12} = 2(\frac16 + \frac{2}{12} + 0) = \frac23$~; par symétrie, toute liaison a $p = 2/3$.
\textbf{17.} Entre ceux des liaisons terminales (0.894) et centrale (0.447) du butadiène.
\textbf{18.} Les six liaisons sont équivalentes par symétrie~; leur longueur, \qty{139.7}{pm}, est comprise
entre celle de la double liaison de l'éthène (\qty{133.9}{pm}) et celle de la liaison simple de l'éthane
(\qty{153.6}{pm}).
\textbf{19.} $0.988|\beta| = \qty{74}{kJ/mol}$, la moitié de celle du benzène~: fermer le cycle double
l'\omterm{def:b2:huckel:delocalisation}{énergie de délocalisation}.
\textbf{20.} $\alpha + 2\beta$, $\alpha + 1.414\beta$ (deux fois), $\alpha$ (deux fois), $\alpha - 1.414\beta$ (deux fois),
$\alpha - 2\beta$.
\textbf{21.} Les deux derniers électrons vont chacun dans une orbitale de la paire non liante~: une couche
ouverte, \omterm{def:b2:huckel:aromatic}{antiaromatique}.
\textbf{22.} Elle se replie en baignoire~: les orbitales p de doubles liaisons voisines ne sont plus
parallèles, les liaisons alternent, et la molécule se comporte comme un polyène.
\textbf{23.} Quatre électrons~: la même paire à moitié remplie, en pire, car la molécule ne peut
éviter la planéité~; elle se déforme en rectangle et est extrêmement réactive.
\textbf{24.} Un système conjugué monocyclique plan est \omterm{def:b2:huckel:aromatic}{aromatique} avec $4N + 2$ électrons $\pi$,
\omterm{def:b2:huckel:aromatic}{antiaromatique} avec $4N$.
\textbf{25.} $\boldsymbol{|\beta| \approx \qty{75}{kJ/mol}}$.
\end{solution}
